\section{Multi-Phi Spectral Aggregation}
\label{sec:multiphi}

% PAGE TARGET: ~7 pages.
% Andrew-style "revised technique" chapter. The default pipeline used
% for all results in Section 4.

The multi-phi detector is the pipeline used to produce every result
in \Cref{sec:benchmark}. It is built on a single observation about
the structure of a DRP stack at fixed elevation: as the source
azimuth $\phi$ rotates around the folio, the laid lines remain in
the same position, but the rest of the image -- paper texture,
fibre clumps, ink edges -- changes with $\phi$. A frequency
estimator that sums over several azimuthal views therefore
accumulates the laid-line signal while averaging the
$\phi$-random noise down. Analyzing mathematically, averaging $N$
independent measurements of a constant signal improves the
signal-to-noise ratio by a factor of $\sqrt{N}$, where $N$ is the
number of azimuthal samples in the DRP stack, typically of order
ten to twenty.
% TODO: Figure 3.1 -- a row of grazing images at four different phi
% values illustrating that the laid lines are visible in the same
% place in all four, while the surrounding texture changes.

\subsection{Collecting one grazing image per phi}
\label{sec:multiphi:collect}

For each azimuth $\phi$ in the DRP stack, the image at the steepest
available elevation $\theta$ is selected. The steepest available
elevation is used because the laid-line contrast is strongest under
grazing illumination, as already noted in
\Cref{sec:first_attempts:simple}. The result is a list of $N$
greyscale images of identical size, one per $\phi$, that share
identical pixel registration with the underlying folio. These images
go through the same background subtraction pipeline to enhance contrast
and relevant features.

\subsection{Aggregating power spectra across azimuth}
\label{sec:multiphi:aggregate}

Each of the $N$ images is rotated so that the laid lines lie
vertically; then, a two-dimensional fast Fourier
transform is taken. The two-dimensional power spectrum is summed
along the vertical axis to give a one-dimensional radial profile
$P_i(f)$ on the horizontal frequency axis. The profile is then
cut to the plausible period band of 8 to 80 pixels. Doing so removes
the low-frequency power from page curvature and multiple harmonics of
the laid lines, and leave the laid line interval estimation to a reasonable range
such that all ground truths of folios in the benchmark are included.

Before the profiles are summed, each one is normalised by
its total band power. This step is essential; the different
azimuthal views of the same folio differ in absolute brightness,
because the light source of the device is not perfectly homogenic across angles, and
because some azimuths catch more reflection from neighbouring
surfaces than others. Without normalisation, the brightest image
would dominate the sum and the advantage of combining multiple azimuthal views would be lost.
After normalisation each image contributes a unit-area spectrum, and
the sum
\begin{equation}
  P_{\mathrm{agg}}(f) \;=\; \sum_{i=1}^{N}
      \frac{P_i(f)}{\int_{\text{band}} P_i(f') \, \mathrm{d}f'}
  \label{eq:agg_power}
\end{equation}
gives the aggregated radial power. The dominant period is estimated as
 the argument of maximum $P_{\mathrm{agg}}$, and the
corresponding frequency $f_{\mathrm{peak}}$ defines the angular
frequency $\omega = 2\pi f_{\mathrm{peak}}$ at which the next
subsection fits the phase.
% TODO: Figure 3.2 -- per-image normalised spectra (thin grey lines)
% overlaid on the aggregated spectrum (thick black line) for one
% folio, showing that the per-image peak is washed out by noise but
% the aggregated peak is sharp.

\subsection{Per-image phase fit and polarity correction}
\label{sec:multiphi:phase}

Knowing the period is not enough to place the grid. The phase of the
laid-line modulation has to be fitted too, and here the multi-phi
data introduces a subtle complication that does not appear in the
single-image case. For each image the one-dimensional
signal $s_i(x)$ is fitted by the cosine and sine components at the
 angular frequency $\omega$, and the per-image phase $\phi_i$
is recovered from the inverse tangent of the two
inner products. This part is identical to the single-image phase
fit of \Cref{sec:first_attempts:simple} except that it is repeated
$N$ times.

The subtle complication is the following. As the illumination
azimuth rotates around the page, the wire shadow that is
characteristic of grazing reflectance flips from one side of each
laid line to the other. Two grazing images $180^{\circ}$ apart in
$\phi$ will cast the shadow of the same wire on opposite sides,
and the fitted per-image phases differ by approximately $\pi$ even
though the underlying period is identical. Simply averaging the
$N$ raw phases would give a meaningless mean phase whose resultant
length is close to zero.

The pipeline solves this by anchoring on the image with the highest
weight at the detected peak, that is, the image that contributes
most to the aggregated spectrum, and then flipping by $\pi$ any
image whose anchor-relative phase difference exceeds $\pi/2$. After
the polarity correction, the resultant length of the weighted
circular mean rises from a typical raw value below $0.3$ to a
post-alignment value above $0.8$ for the folios in the benchmark.
The post-alignment resultant length is reported alongside the
final density as a diagnostic of how coherent the multi-phi
evidence was, and a low value flags the case where the laid-line
signal was so weak that the polarity-flip detection itself failed.
% TODO: Figure 3.3 -- per-image phase polar plot before and after
% polarity correction, with marker size proportional to per-image
% weight. The "before" panel shows phases scattered around two
% antipodal directions; the "after" panel shows them clustered.

\subsection{Weighted mean phase and half-period auto-correction}
\label{sec:multiphi:final}

The aligned per-image phases are then combined into a single mean
phase by the amplitude-weighted circular mean, with the weights
taken from the per-image power at the detected peak. The mean
phase together with the locked period defines the grid positions in
the rotated frame of the representative image. The representative
image is taken to be the highest-weight image and is the one used
later for the wire-width estimate and the overlay shown to the
reader.

There remains a final ambiguity. The polarity correction of the
previous subsection forces all per-image phases to agree, but it
does not by itself fix the absolute polarity of the resulting
combined phase. The combined grid can therefore land either on the
dark wire-shadow side or on the bright inter-wire side of the
broadband signal. The pipeline resolves this by sampling the
column-mean of the representative image at the grid positions and
at positions shifted by half a period. If the grid positions have
higher intensity than the half-period-shifted positions, the grid
is sitting on the bright side and the phase is shifted by $\pi$ so
that the final grid always lands on the wire shadows. The shift
makes the overlay independent of which side the polarity-anchor
image happened to favour and gives a more physically meaningful
output, since codicologists describe densities by reference to the
wire positions.

The output of the multi-phi detector therefore comprises the
dominant period, the absolute phase, the line direction, and the
resulting grid in the original-image coordinate frame, along with
the diagnostic resultant lengths and per-image weights described
above.
% TODO: Figure 3.4 -- a representative overlay showing the final
% grid drawn on the bg-subtracted grazing image used as the
% reference. Choose Kk1-5 f5v as the example since this is the
% best-case folio in Section 4.

\subsection{Evaluation metrics}
\label{sec:multiphi:evaluation}

A detector that returns only a density figure is hard to deploy at
scale. The figure may be right or it may be wrong, and a downstream
codicologist who is presented with a thousand folio-level numbers
has no way to know which to trust. A single goodness-of-fit score
is not enough either, because a laid-line detection can fail in
several qualitatively different ways: the spectral peak can be
correct in location but not in shape; the spatial regularity of the
wires can break down; the direction in which the signal is strongest
can be only marginally better than a random direction; the
detection can ride on a single lucky azimuth and collapse when that
azimuth is removed; or the wires can be uniform on average but vary
locally across the page. The pipeline therefore reports five
quantities, each of which is sensitive to one of these failure
modes and largely insensitive to the others. All five are derived
from the detector output itself and require no external ground
truth, so they can flag a low-confidence case before any manual
check is done.

The harmonic fit quality $R^2$ asks whether the spectral peak is a
real harmonic structure. A Fourier series with the locked
fundamental and the first three harmonics is fitted to the broadband
signal, and the coefficient of determination is reported. A value
close to one means that almost all of the variance of the signal is
explained by harmonics of the detected period, which is what a real
laid-line pattern produces. A low value flags a peak that was the
largest in the spectrum but did not correspond to a clean periodic
structure, for example a peak caused by paper curvature or by a
single high-amplitude artefact.

The empirical interval distribution asks the same question in the
spatial domain. Peaks of the broadband signal are detected within a
narrow tolerance of the locked period, and the distribution of
peak-to-peak distances is reported as a median, an interquartile
range, and a coefficient of variation. The median gives the
per-folio density once the optical scale is applied. The
coefficient of variation gives a robust measure of how regular the
laid lines actually are; a value below about ten per cent indicates
a well-formed mould signature, and a much larger value either
indicates physical damage to the folio or a detection that has
attached to noise rather than to a real periodicity.
% TODO: Figure 3.5 -- histogram of laid-line intervals on Kk1-5 f5v,
% with the locked period marked by a vertical line.

The self-consistency contrast asks whether the detected
direction itself is special, by comparing the detector score along
the inferred laid-line orientation with its scores along a sequence
of randomly chosen non-laid-line directions on the same image. The
result is reported as the $z$-score of the laid-line score against
this empirical null distribution. A value above about five indicates
that the detected direction is uniquely periodic, which is the
expected behaviour for a folio with real laid lines. A value near
zero indicates that the detector would have reported a similar score
for a random direction, which is the signature of a folio that does
not contain any usable laid-line signal at all.

The split-half stability asks whether the detection depends on any
one image. The $N$ grazing images are partitioned at random into
two disjoint halves; the detector is run on each half independently,
and the standard deviation of the recovered period over many such
partitions is reported. A small standard deviation, of the order of
one per cent of the locked period, indicates that the same period
would have been recovered from a quite different subset of the data.
A large standard deviation indicates that the detection rests on
one or two azimuthal views and would have changed substantially had
those views been absent.

The per-segment wire-width distribution asks whether the laid wires
behave consistently across the page. The folio is divided into a
small number of stripes along the laid-line direction, and the wire
full-width-at-half-maximum is fitted independently within each
stripe. The interquartile range of the per-stripe widths captures
spatial variation that the global wire-width figure smooths over,
and a large range flags either local damage, surface curvature, or
a sheet that was not properly flattened during imaging.

The five metrics overlap deliberately. A single failure mode
typically degrades at least two of them at once, so a folio that is
healthy on all five is more trustworthy than one that is healthy on
any single metric; a folio that is bad on one but healthy on the
other four can usually be diagnosed by reading which metric is bad
and locating the relevant feature in the corresponding plot.
\Cref{sec:benchmark} reports the five for every folio in the
benchmark, and the failure cases collected in
\Cref{sec:benchmark:failures} are organised around which metric or
combination of metrics first identified each problem.
